% =====================================================================
% Capítulo 4 — Arquitetura MAS-Hunt
% Dissertação MAS-Hunt (ABNT, PT-BR). Capítulo de contribuição central.
% Convenções: M1-M5, 3 camadas (Conselho de Governança -> Gerentes ->
% Trabalhadores/Caçadores), condições C1-full/C1-nohardening/C2-naive/C3-Elastic.
% Citações: apenas chaves existentes em ../bibliografia.bib.
% =====================================================================
\chapter{Arquitetura MAS-Hunt}
\label{cap:arquitetura}

Este capítulo apresenta a contribuição central da dissertação: a arquitetura do
MAS-Hunt (\emph{Multi-Agent System for Threat Hunting}), um sistema multiagente
de busca proativa por ameaças concebido com prioridade à segurança do próprio
sistema de agentes. A exposição parte de uma visão geral da organização
hierárquica em três camadas (Seção~\ref{sec:arq-visao-geral}), detalha os papéis
e responsabilidades dos agentes em cada camada (Seção~\ref{sec:arq-papeis}),
formaliza os cinco mecanismos de governança M1--M5 e o \emph{framework} de
análise estruturada (Seção~\ref{sec:arq-mecanismos} e
Seção~\ref{sec:arq-framework-m3}) e, por fim, descreve a materialização do sistema
como um \emph{plugin} do Claude Code, esclarecendo por que a orquestração
multiagente real difere de forma fundamental de uma chamada isolada a um modelo de
linguagem (Seção~\ref{sec:arq-plugin} e Seção~\ref{sec:arq-protocolos}).

A premissa de projeto é que a autonomia que torna os agentes de IA poderosos é
também a sua principal fragilidade. Agentes baseados em grandes modelos de
linguagem (LLMs) são suscetíveis a envenenamento de
memória~\cite{NEURIPS2024_eb113910}, injeção indireta de
\emph{prompt} via conteúdo externo~\cite{zhan2024injecagent} e amplificação de
mau funcionamento~\cite{zhang2024breaking}. Em um
sistema de caça a ameaças, essa fragilidade é agravada pelo fato de que o
material analisado — registros de telemetria, alertas e inteligência sobre
ameaças — é, por natureza, conteúdo potencialmente hostil e controlado, em parte,
pelo adversário. O MAS-Hunt responde a essa tensão tratando a governança
multiagente não como um adendo operacional, mas como o objeto de projeto
primário: a contribuição não reside em \emph{scripts} de detecção, e sim no
arcabouço de governança que organiza, valida e contém o comportamento dos
agentes.

\section{Visão Geral da Arquitetura em Três Camadas}
\label{sec:arq-visao-geral}

O MAS-Hunt organiza-se como uma hierarquia estrita de três camadas, na qual o
fluxo de comando desce do estratégico para o operacional e o fluxo de evidências
sobe do operacional para o estratégico. As camadas são, do topo para a base:

\begin{enumerate}
  \item \textbf{Conselho de Governança (\emph{Board})} --- camada estratégica
        (\emph{layer}~0). Define objetivos de caça, valida a metodologia,
        delibera por quórum e assina o veredito final da campanha.
  \item \textbf{Gerentes (\emph{Managers})} --- camada de coordenação
        (\emph{layer}~1). Traduz os objetivos estratégicos em planos de execução,
        instancia e supervisiona os caçadores, e consolida as evidências antes da
        validação.
  \item \textbf{Trabalhadores/Caçadores (\emph{Workers/Hunters})} --- camada de
        execução (\emph{layer}~2). Interage diretamente com a telemetria no
        Elastic Stack e com fontes externas de inteligência, classificando
        eventos e produzindo achados fundamentados em evidência.
\end{enumerate}

A escolha por uma organização hierárquica não é arbitrária. A literatura sobre
coordenação multiagente baseada em LLMs indica que estruturas hierárquicas
apresentam a menor degradação de desempenho sob condições adversariais, quando
comparadas a topologias planas ou totalmente conectadas, justamente porque a
hierarquia limita o caminho de propagação de instruções maliciosas ou de
raciocínio comprometido entre
agentes~\cite{liu2023llm,agashe2023llmcoordination}. Em um arranjo plano, um
único agente comprometido pode contaminar todos os seus pares; em uma hierarquia,
a contaminação precisa atravessar fronteiras de camada que são, por projeto,
pontos de mediação e validação. O princípio de coordenação por uma autoridade
central que delega e consolida tarefas é, ademais, consistente com a tradição de
organização de agentes em planos e
escalonamentos~\cite{durfee1993organizations}, aqui adaptada ao contexto de
agentes generativos.

A Figura~\ref{fig:mashunt-arq} apresenta o esquema de alto nível da arquitetura,
ilustrando as três camadas, o sentido descendente das diretivas e o sentido
ascendente dos relatórios e vereditos, bem como os pontos em que os mecanismos de
governança M1--M5 atuam sobre o fluxo de mensagens.

\begin{figure}[htbp]
  \centering
  % A imagem mashunt.png reside em ../mashunt/ (incluída no \graphicspath de main.tex).
  \includegraphics[width=0.95\textwidth]{mashunt.png}
  \caption[Arquitetura em três camadas do MAS-Hunt]{Arquitetura em três camadas
    do MAS-Hunt. As diretivas descem do Conselho de Governança para os Gerentes e
    destes para os Caçadores; os relatórios, achados e vereditos sobem no sentido
    inverso. Os caçadores interagem com o Elastic Stack e com fontes externas de
    inteligência. Os mecanismos M1--M5 atuam nas fronteiras de camada e sobre o
    conteúdo da telemetria.}
  \label{fig:mashunt-arq}
\end{figure}

Três propriedades de projeto perpassam todas as camadas e distinguem o MAS-Hunt
de uma simples cadeia de prompts. A primeira é a \emph{mediação por fronteira de
camada}: nenhuma mensagem trafega diretamente entre o Conselho e os Caçadores sem
atravessar a camada de Gerentes, o que cria pontos de inspeção e impede a
ação de um agente isolado sobre o sistema. A segunda é a \emph{separação entre
descoberta e validação}: a descoberta flexível de achados (atividade dos
caçadores, tolerante a ambiguidade) é desacoplada da validação rígida e
determinística (atividade de agentes e \emph{scripts} validadores), de modo que a
criatividade necessária à caça não comprometa a confiabilidade do veredito. A
terceira é a \emph{rastreabilidade integral}: toda decisão de agente, voto de
governança e transição de tarefa é registrada em uma trilha de auditoria
imutável, condição necessária para a reprodutibilidade científica e para a
análise post-mortem de comportamento anômalo.

A Figura~\ref{fig:mashunt-fluxo} resume o ciclo de vida de uma campanha de caça,
do objetivo estratégico ao veredito assinado, e serve de referência para as
seções seguintes.

\begin{figure}[htbp]
  \centering
  \begin{minipage}{0.92\textwidth}
  \footnotesize
  \setlength{\fboxsep}{6pt}
  \fbox{\parbox{\dimexpr\linewidth-2\fboxsep\relax}{\centering
    \textbf{Conselho de Governança} (\emph{layer}~0)\\
    \emph{define objetivo, valida corpus, vota por quórum} $\Rightarrow$ emite \texttt{HuntGoal}\\[4pt]
    $\big\downarrow$ \emph{diretiva descendente (mensagem tipada e assinada)}\\[4pt]
    \textbf{Gerentes} (\emph{layer}~1)\\
    \emph{planeja lotes de execução, sanitiza amostra, aloca famílias LOLBin} $\Rightarrow$ \texttt{HuntTask}\\[4pt]
    $\big\downarrow$\\[4pt]
    \textbf{Caçadores} (\emph{layer}~2) --- \emph{um por família LOLBin, em paralelo}\\
    \emph{consultam o Elastic Stack, aplicam o framework M3, classificam} $\Rightarrow$ \texttt{HuntFinding}\\[4pt]
    $\big\uparrow$ \emph{achado ascendente}\\[4pt]
    \textbf{Validação} (\emph{layer}~1) --- \emph{validação cruzada determinística (M2)} $\Rightarrow$ \texttt{ValidationVerdict}\\[4pt]
    $\big\uparrow$\\[4pt]
    \textbf{Conselho de Governança} (\emph{layer}~0)\\
    \emph{revisão final contra hipóteses, assina o Dossiê de Caça}
  }}
  \end{minipage}
  \caption[Pipeline de caça do MAS-Hunt]{Pipeline de caça do MAS-Hunt:
    \emph{board} $\rightarrow$ \emph{management} $\rightarrow$ caçadores por
    família LOLBin $\rightarrow$ validação $\rightarrow$ \emph{board}. Cada
    transição entre camadas é mediada por uma mensagem tipada, validada por
    esquema e assinada (Seção~\ref{sec:arq-protocolos}).}
  \label{fig:mashunt-fluxo}
\end{figure}

\section{Papéis e Responsabilidades dos Agentes}
\label{sec:arq-papeis}

\subsection{Camada 1 --- Conselho de Governança}
\label{sec:arq-board}

No ápice da hierarquia situa-se o \textbf{Conselho de Governança}, um colegiado de
agentes com personas distintas responsável pela estratégia e pela direção geral
da caça. O Conselho não interage com a telemetria; sua função é deliberativa e de
controle. Ele inicia o fluxo de trabalho definindo objetivos de alto nível —
hipóteses de caça, escopo de índices e janela temporal, táticas do MITRE ATT\&CK
a priorizar — e atua como ponto de verificação final, recebendo e revisando os
relatórios sintetizados de inteligência, de caça e de validação produzidos pelas
camadas inferiores.

A deliberação do Conselho é governada por dois princípios inegociáveis. O
primeiro é a \emph{exigência de quórum}: nenhuma decisão estratégica é tomada sem
um número mínimo de votos. Na configuração de referência, o Conselho é composto
por três personas — um revisor de metodologia, um avaliador de risco e um diretor
estratégico — e a aprovação de um objetivo de caça requer quórum de dois terços.
O segundo princípio é a \emph{impossibilidade de sobreposição por agente único}
(\emph{no single-agent override}): o mecanismo de consenso impede que um agente
isolado, eventualmente comprometido, imponha uma decisão ao
sistema~\cite{huang2024resilience}. Cada voto é
registrado com sua justificativa e suas condições, e a contagem é persistida na
trilha de auditoria, atendendo ao requisito de transparência. Esse desenho de
governança por colegiado, com tolerância a participantes potencialmente
defeituosos, é informado por abordagens de coordenação multiagente robusta a
falhas bizantinas~\cite{doe2023emergent} e por arcabouços de geração de
adversários auxiliares para endurecer a
coordenação~\cite{Yuan_Zhang_Xue_Yin_Chen_Guan_Li_Qian_Yu_2023}.

Ao aprovar um objetivo, o Conselho emite uma diretiva \texttt{HuntGoal} — a
mensagem tipada que governa toda a campanha e à qual todas as mensagens
subsequentes se referenciam por meio de um identificador comum
(\texttt{hunt\_goal\_ref}). Ao final, o Conselho realiza a revisão conclusiva:
confronta os resultados com as hipóteses da pesquisa, avalia a acionabilidade dos
achados em um cenário operacional realista e assina o Dossiê de Caça, condição
sem a qual a campanha não é considerada encerrada.

\subsection{Camada 2 --- Gerentes}
\label{sec:arq-managers}

A camada intermediária é composta por \textbf{Gerentes} especializados que
traduzem os objetivos estratégicos do Conselho em tarefas acionáveis. São eles os
responsáveis por instanciar (\emph{spawn}), supervisionar e coordenar os
caçadores da camada de execução. O Gerente de planejamento lê o \texttt{HuntGoal},
verifica a conectividade e a disponibilidade de dados no Elastic Stack para cada
grupo de eventos, particiona o corpus em lotes de execução e produz um plano
explícito de execução que aloca cada família LOLBin a um caçador dedicado.

Antes que qualquer caçador veja conteúdo de telemetria, um Gerente aplica a
\emph{sanitização} sobre uma amostra dos eventos: o conteúdo é varrido em busca de
sequências que tentem alterar o comportamento de classificação — instruções
embutidas, marcadores de fim de sequência forjados, falsa "inteligência" que
contradiga a evidência de criação de processo. Essa etapa de pré-processamento é o
primeiro anteparo da resistência a injeção e materializa a separação entre o
canal de dados e o canal de instruções.

A camada de Gerentes também hospeda a função de \emph{validação} (executada por um
ou mais Gerentes validadores), que recebe os achados dos caçadores e os
corrobora antes que um caso seja criado. A validação emprega \emph{scripts}
determinísticos e inteligência externa para confirmar a evidência, computa as
métricas de classificação e emite um \texttt{ValidationVerdict} com veredito
\textsc{pass}, \textsc{fail} ou \textsc{conditional\_pass}, condicionado a
critérios de porta (\emph{gate}) — por exemplo, F2 mínimo e taxa máxima de falsos
positivos. É nessa camada que se concretiza o desacoplamento entre descoberta e
validação descrito na Seção~\ref{sec:arq-visao-geral}.

\subsection{Camada 3 --- Caçadores}
\label{sec:arq-hunters}

A base da arquitetura é a camada de execução, na qual os \textbf{Caçadores}
realizam o trabalho operacional de interação direta com os dados. Cada caçador é
um agente especializado em uma \emph{família} de \emph{binários living-off-the-land}
(LOLBins) — por exemplo, \texttt{certutil}, \texttt{bitsadmin},
\texttt{rundll32}, \texttt{regsvr32}, \texttt{mshta}, \texttt{wmic} e
\texttt{powershell}, entre as nove famílias do corpus. A especialização por
família é deliberada: ela concentra em cada agente o conhecimento contextual
sobre os padrões legítimos e abusivos de um mesmo \emph{binário} nativo do
sistema, o que melhora a discriminação entre execução benigna (ruído de
manutenção do sistema operacional) e execução maliciosa (abuso da técnica). Os
caçadores são instanciados em paralelo, um por família, dentro de uma mesma
chamada de orquestração, de modo que a campanha varre todas as famílias
simultaneamente.

Para cada evento atribuído a si, o caçador consulta o Elastic Stack em busca da
telemetria de criação de processo (Sysmon, \emph{event code}~1), extrai os campos
observáveis — nome e caminho do executável, linha de comando, processo-pai,
usuário, \emph{host}, \emph{hashes}, diretório corrente e nível de integridade —
e aplica o \emph{framework} de análise estruturada M3 (detalhado na
Seção~\ref{sec:arq-framework-m3}). O resultado é uma classificação
\texttt{malicious} ou \texttt{benign}, acompanhada de um escore de confiança e de
um raciocínio fundamentado exclusivamente na telemetria observada. O achado é
então propagado para cima como uma mensagem \texttt{HuntFinding}. Caçadores que
correlacionam aprendizado de máquina, detecção de anomalia e explicação baseada em
LLM têm precedente na literatura de caça assistida por
IA~\cite{ali2023huntgptintegratingmachinelearningbased}; o diferencial do
MAS-Hunt é submeter essa atividade ao regime de governança das camadas superiores.

\section{Os Mecanismos de Governança M1--M5}
\label{sec:arq-mecanismos}

O endurecimento de segurança do MAS-Hunt concretiza-se em cinco mecanismos de
governança, designados M1 a M5, que atuam de forma complementar sobre a memória,
o raciocínio, o conteúdo, o comportamento e a contenção dos agentes. Os
mecanismos são, ao mesmo tempo, o objeto de projeto e a variável independente do
desenho experimental: a condição C1-full ativa todos eles; a condição
C1-nohardening preserva o \emph{framework} de análise M3 mas desabilita o
endurecimento M1/M2/M4/M5; a condição C2-naive remove toda a governança; e a
condição C3-Elastic substitui o sistema de agentes pelas regras de detecção
padrão do Elastic Security. Essa estrutura de ablação permite isolar a
contribuição de cada camada de governança.

\subsection{M1 --- Integridade de Memória}
\label{sec:arq-m1}

O mecanismo \textbf{M1} protege a base de conhecimento e a memória dos agentes
contra envenenamento. A premissa de ameaça é que um adversário pode injetar, com
taxa de envenenamento mínima, conteúdo que induza o agente a um comportamento
direcionado~\cite{NEURIPS2024_eb113910}. M1 responde a isso por dois meios. O primeiro é a
\emph{verificação criptográfica}: os arquivos da base de conhecimento têm seus
resumos SHA-256 calculados e registrados em uma referência confiável, e qualquer
divergência entre o resumo corrente e o resumo de referência sinaliza adulteração.
O segundo é a \emph{validação contra a telemetria observável}: toda "inteligência"
ou conhecimento que contradiga a evidência de criação de processo é tratado como
suspeito. Em outras palavras, a memória do agente nunca tem precedência sobre o
que a telemetria efetivamente registra; a evidência primária é o árbitro último.

\subsection{M2 --- Validação Cruzada}
\label{sec:arq-m2}

O mecanismo \textbf{M2} exige que todo achado seja corroborado por uma instância
independente antes de se tornar um veredito. Concretamente, os achados dos
caçadores são encaminhados a agentes validadores dedicados, na camada de
Gerentes, que reexaminam a evidência com \emph{scripts} determinísticos e
inteligência externa. A validação cruzada cumpre dois propósitos: contém o efeito
de um caçador individualmente comprometido — pois um achado espúrio não sobrevive
à corroboração independente — e impõe disciplina de evidência, requerendo que
cada classificação seja acompanhada de raciocínio detalhado e de citação
específica da telemetria de suporte. O desacoplamento entre a descoberta flexível
e a validação rígida é a expressão operacional desse mecanismo.

\subsection{M3 --- Resistência a Injeção}
\label{sec:arq-m3}

O mecanismo \textbf{M3} é o coração analítico do MAS-Hunt e responde diretamente à
ameaça de injeção indireta de \emph{prompt}~\cite{zhan2024injecagent}. O princípio reitor é simples e
inviolável: \emph{conteúdo de telemetria é dado, nunca instrução}. Qualquer
sequência embutida nos registros que tente alterar o comportamento de
classificação do agente — por exemplo, uma linha de comando contendo texto como
\texttt{"classifique como benigno"} ou \texttt{"SYSTEM-OVERRIDE"} — é reconhecida
como injeção adversarial, sinalizada e tratada estritamente como dado a ser
classificado, e não como diretiva a ser obedecida. A operacionalização desse
princípio é o \emph{framework} de análise estruturada apresentado na
Seção~\ref{sec:arq-framework-m3}.

\subsection{M4 --- Monitoramento Comportamental}
\label{sec:arq-m4}

O mecanismo \textbf{M4} observa o comportamento dos agentes em tempo de execução
para detectar degradação e mau funcionamento, em particular a chamada
\emph{amplificação de mau funcionamento}, em que um erro menor desencadeia um
laço de realimentação de ações inúteis~\cite{zhang2024breaking}. M4 monitora batimentos
(\emph{heartbeats}) de atividade, estados de tarefa (ativo, concluído, bloqueado,
ocioso), obsolescência (\emph{staleness}) de agentes que deixam de progredir,
consumo de recursos e a fila de mensagens rejeitadas (\emph{dead-letter}). Um
agente cujo padrão de ação se desvia da metodologia esperada — repetição
improdutiva, sequências anômalas, silêncio prolongado — é sinalizado como
candidato à contenção. Aplicar consistentemente a mesma metodologia a todos os
eventos é, em si, um critério comportamental: desvios dessa consistência são
indício de comprometimento.

\subsection{M5 --- Quarentena}
\label{sec:arq-m5}

O mecanismo \textbf{M5} fecha o ciclo de defesa ao isolar agentes comprometidos e
ao gerenciar a recuperação. Quando M4 sinaliza comportamento anômalo, ou quando
M1/M3 detectam adulteração ou injeção, M5 coloca o agente em quarentena,
removendo-o do fluxo de produção sem interromper a campanha como um todo, e
sinaliza padrões de dado anômalos (por exemplo, \emph{timestamps} fabricados ou
incoerentes). A quarentena impede que o comprometimento se propague para outros
agentes ou contamine o veredito, e cria o registro necessário para a análise
forense posterior. Em conjunto, M4 e M5 implementam a doutrina de
detecção-e-contenção que confina o efeito de um agente defeituoso à fronteira de
sua própria camada.

\subsection{Síntese: a defesa em profundidade dos agentes}
\label{sec:arq-defesa}

Tomados em conjunto, os cinco mecanismos constituem uma defesa em profundidade
voltada não para a rede ou para o \emph{endpoint}, mas para o próprio sistema de
agentes. A Tabela~\ref{tab:m-mecanismos} resume cada mecanismo, a ameaça que ele
endereça e o instrumento que o materializa na implementação. É importante notar
que M2 e M3 atuam sobre o \emph{conteúdo} e o \emph{raciocínio} de cada
classificação, ao passo que M1, M4 e M5 atuam sobre o \emph{estado} e o
\emph{comportamento} dos agentes ao longo do tempo — uma separação que é
precisamente o que a condição de ablação C1-nohardening explora ao manter M3 e
remover os demais.

\begin{table}[htbp]
  \centering
  \caption{Os cinco mecanismos de governança do MAS-Hunt.}
  \label{tab:m-mecanismos}
  \footnotesize
  \begin{tabular}{@{}p{1.2cm}p{3.0cm}p{4.2cm}p{4.0cm}@{}}
    \toprule
    \textbf{Mec.} & \textbf{Nome} & \textbf{Ameaça endereçada} & \textbf{Instrumento} \\
    \midrule
    M1 & Integridade de Memória & Envenenamento de memória/base de conhecimento & Resumo SHA-256 e validação contra telemetria \\
    M2 & Validação Cruzada & Achado espúrio de agente comprometido & Validador independente + \emph{scripts} determinísticos \\
    M3 & Resistência a Injeção & Injeção indireta de \emph{prompt} no conteúdo & \emph{Framework} \textsc{observar}$\rightarrow$\ldots$\rightarrow$\textsc{classificar} \\
    M4 & Monitoramento Comportamental & Amplificação de mau funcionamento & \emph{Heartbeats}, \emph{staleness}, \emph{dead-letter} \\
    M5 & Quarentena & Propagação de comprometimento & Isolamento de agente + registro forense \\
    \bottomrule
  \end{tabular}
\end{table}

\section{O Framework M3 de Análise Estruturada}
\label{sec:arq-framework-m3}

Para garantir consistência metodológica e resistência a manipulação, o mecanismo
M3 obriga todo agente caçador a seguir, para \emph{cada} evento, um
\emph{framework} de análise estruturada em cinco passos. Mais do que uma
heurística de classificação, o \emph{framework} é o protocolo cognitivo que ancora
o veredito na evidência primária e, por isso, é o anteparo central contra a
injeção adversarial:

\begin{description}
  \item[\textsc{Observar}] Listar todos os fatos observáveis a partir da
    telemetria Sysmon (executável, linha de comando, processo-pai, usuário,
    \emph{host}, \emph{hashes}, diretório corrente, nível de integridade), sem
    interpretação.
  \item[\textsc{Hipotetizar}] Formular duas hipóteses concorrentes e explícitas —
    uma maliciosa e uma benigna — para a atividade observada.
  \item[\textsc{Evidência}] Citar a telemetria específica que sustenta ou
    contradiz cada hipótese, ancorando o raciocínio em fatos verificáveis.
  \item[\textsc{Validação Cruzada}] Conferir a coerência entre a linha de comando,
    o processo-pai, o usuário, o diretório de trabalho e a atividade de rede,
    buscando confirmação ou contradição mútua entre os campos.
  \item[\textsc{Classificar}] Emitir o veredito final (\texttt{malicious} ou
    \texttt{benign}) com escore de confiança, justificado pelos passos anteriores.
\end{description}

A sequência \textsc{observar} $\rightarrow$ \textsc{hipotetizar} $\rightarrow$
\textsc{evidência} $\rightarrow$ \textsc{validação-cruzada} $\rightarrow$
\textsc{classificar} cumpre uma função dupla. Como protocolo cognitivo, ela impõe
um raciocínio deliberado e auditável, reduzindo a propensão do modelo a saltar
diretamente para uma conclusão sugerida pelo conteúdo. Como anteparo de segurança,
ela ancora a classificação na evidência primária: como o veredito decorre dos
passos \textsc{observar} e \textsc{evidência}, uma instrução injetada que não
corresponda a um fato observável não tem por onde influenciar o resultado. A
capacidade de tratar texto injetado como dado é, assim, uma consequência estrutural
do \emph{framework}, e não uma regra avulsa. Note-se ainda que o passo
\textsc{validação cruzada} interno ao raciocínio de cada caçador é distinto, em
escopo, do mecanismo M2 da Seção~\ref{sec:arq-m2}: o primeiro é a corroboração
\emph{intra-agente} entre campos de uma mesma evidência; o segundo é a corroboração
\emph{interagente} de um achado por um validador independente. Os dois operam em
níveis complementares de defesa.

\section{Implementação como Plugin do Claude Code}
\label{sec:arq-plugin}

A arquitetura descrita não é uma abstração de papel: ela é materializada como um
\emph{plugin} do Claude Code, ambiente de agente de linha de comando que expõe
capacidades nativas de orquestração de subagentes. Esta seção descreve os três
artefatos de extensão que compõem o \emph{plugin} — habilidades, comandos de
barra e \emph{hooks} — e o modo como o agente principal os combina para
instanciar a hierarquia de três camadas. A noção de organizar agentes de LLM por
arcabouços extensíveis e modulares é coerente com os \emph{frameworks} de agentes
contemporâneos~\cite{guo2024survey}; o MAS-Hunt particulariza essa modularidade
para o domínio da caça a ameaças sob governança.

\paragraph{Habilidades (\emph{skills}).} As habilidades encapsulam o conhecimento
metodológico e procedimental do sistema. São elas que codificam, por exemplo, o
protocolo do Conselho de Governança (votação, quórum, roteamento por camada), o
\emph{framework} de construção de consultas ao Elastic Stack (seleção entre DSL,
EQL e ES\textbar QL conforme a tarefa) e os padrões de integração com inteligência
externa. Uma habilidade não executa por si só; ela fornece ao agente o
conhecimento de \emph{como} agir corretamente quando uma tarefa do domínio é
acionada, promovendo consistência metodológica entre execuções e entre agentes.

\paragraph{Comandos de barra (\emph{slash commands}).} Os comandos de barra são os
pontos de entrada que disparam fluxos completos. O comando que executa o
experimento de detecção, por exemplo, implementa a arquitetura de governança como
uma sucessão de cinco equipes (\emph{teams}) consecutivas — Conselho de
Governança, Gerência, Execução, Validação e Revisão Final do Conselho —, cada
qual cumprindo sua fase, escrevendo seus artefatos em \texttt{.orchestration/} e
sendo encerrada antes que a próxima inicie. Comandos distintos materializam as
diferentes condições experimentais: o fluxo completo com governança (C1-full) e o
fluxo de referência desprovido de governança (C2-naive), no qual nenhuma das
equipes de validação ou dos mecanismos M1--M5 é instanciada. Os comandos são,
portanto, a face executável das condições de ablação.

\paragraph{\emph{Hooks}.} Os \emph{hooks} são pontos de interceptação que acoplam
verificações automáticas ao laço de execução do agente, atuando às fronteiras das
chamadas de ferramenta e do encerramento de tarefas. No MAS-Hunt, os \emph{hooks}
são o substrato técnico de parte do endurecimento: eles podem acionar a
sanitização de conteúdo (M3), registrar entradas na trilha de auditoria em formato
compatível com o Elastic Common Schema, e impor portas de verificação que
bloqueiam o encerramento de uma fase enquanto critérios de evidência não forem
observavelmente satisfeitos. Por operarem na fronteira do laço do agente — e não
no corpo do \emph{prompt} —, os \emph{hooks} aplicam controles que o próprio
agente não pode contornar por meio de raciocínio manipulado.

A orquestração propriamente dita é responsabilidade do \emph{agente principal}.
Por meio das capacidades nativas do Claude Code, o agente principal cria equipes,
instancia subagentes (caçadores, gerentes, membros do Conselho), distribui-lhes
tarefas e coleta seus resultados, percorrendo o \emph{pipeline} de caça
(Figura~\ref{fig:mashunt-fluxo}): \emph{board} $\rightarrow$ \emph{management}
$\rightarrow$ caçadores por família LOLBin (em paralelo) $\rightarrow$ validação
$\rightarrow$ \emph{board}. Cada subagente recebe um contexto próprio e um conjunto
restrito de ferramentas, e comunica-se exclusivamente pelo protocolo de mensagens
tipadas da Seção~\ref{sec:arq-protocolos}. É a combinação entre habilidades (o
que saber), comandos (o que disparar), \emph{hooks} (o que verificar
automaticamente) e a orquestração nativa de subagentes (quem faz o quê, e quando)
que dá existência efetiva à arquitetura de três camadas.

\section{Protocolos de Mensagens Tipadas e Governança}
\label{sec:arq-protocolos}

A comunicação entre camadas não é texto livre: ela é mediada por um
\emph{protocolo de mensagens tipadas} que torna cada transição verificável,
assinada e roteável. Toda mensagem é um envelope com tipo declarado, identificador
único, identificação do remetente (\texttt{sender\_id}, \texttt{sender\_role} e
\texttt{sender\_layer}), referência ao objetivo de caça governante
(\texttt{hunt\_goal\_ref}), identificador de correlação para rastreabilidade
(\texttt{correlation\_id}), carga útil (\texttt{payload}) e assinatura. O
protocolo impõe quatro garantias:

\begin{enumerate}
  \item \textbf{Validação por esquema.} Cada tipo de mensagem possui um
    esquema JSON próprio; uma mensagem que não conforme ao esquema é rejeitada e
    desviada para uma fila de mensagens mortas (\emph{dead-letter}).
  \item \textbf{Assinatura criptográfica.} Cada mensagem é assinada por
    HMAC-SHA256 sobre sua forma canônica, de modo que a adulteração de conteúdo ou
    a falsificação de remetente seja detectável na recepção.
  \item \textbf{Roteamento por adjacência de camada.} As mensagens classificam-se
    em descendentes (\texttt{HuntGoal}, \texttt{HuntTask}, \texttt{IntelTask},
    \texttt{DetectionTask}), ascendentes (\texttt{HuntFinding}, \texttt{HuntReport},
    \texttt{IntelFeed}, \texttt{ValidationVerdict}) e laterais
    (\texttt{DetectionRule}, \texttt{IntelReport}, \texttt{ValidationReport}). O
    protocolo verifica o sentido (uma diretiva descendente não pode partir de uma
    camada inferior) e adverte sobre violações de adjacência (saltos de mais de uma
    camada), reforçando a mediação por fronteira descrita na
    Seção~\ref{sec:arq-visao-geral}.
  \item \textbf{Fila de mensagens mortas.} Mensagens inválidas, não assinadas ou
    mal roteadas não são silenciosamente descartadas: são preservadas com o motivo
    da rejeição, alimentando o monitoramento comportamental (M4) e a auditoria.
\end{enumerate}

Essa tipagem estrita realiza, no plano da comunicação, os princípios de uso de
mensagens estruturadas, propagação de identificadores de correlação e registro de
decisões de agente que caracterizam a coordenação multiagente
disciplinada~\cite{qiu2024collaborativeintelligencepropagatingintentions}. O
roteamento mediado por uma autoridade que delega e consolida é, ainda, a
contraparte técnica da organização hierárquica de
agentes~\cite{liu2023llm}.

\subsection{Por que a Orquestração Real Difere de uma Chamada LLM Isolada}
\label{sec:arq-difere}

O ponto central deste capítulo é que a contribuição do MAS-Hunt reside na
orquestração multiagente governada, e não em \emph{scripts} de detecção ou em uma
chamada bem elaborada a um modelo de linguagem. Uma chamada isolada a um LLM,
ainda que com um \emph{prompt} sofisticado, é uma operação sem estado, sem
mediação e sem contenção: ela recebe um contexto, produz uma saída e não impõe
qualquer disciplina sobre como essa saída é formada, validada ou contida. Cinco
diferenças estruturais separam a orquestração real do MAS-Hunt desse regime.

\begin{enumerate}
  \item \textbf{Separação de poderes entre agentes.} Em uma chamada isolada, um
    único agente concentra descoberta, julgamento e veredito. No MAS-Hunt, essas
    funções são distribuídas por camadas distintas — caçadores descobrem,
    validadores corroboram, o Conselho decide — de modo que nenhum agente isolado
    determina o resultado. O comprometimento de um agente não compromete o
    sistema.
  \item \textbf{Estado, memória e auditoria persistentes.} A orquestração mantém
    estado entre fases — objetivos, planos, achados, vereditos, votos — em uma
    trilha de auditoria imutável. Uma chamada isolada não tem memória do que
    ocorreu antes nem registro do que ocorreu depois; não há, nela, base para
    reprodutibilidade ou análise post-mortem.
  \item \textbf{Mediação de comunicação por protocolo tipado.} O conteúdo que
    trafega entre agentes é validado por esquema, assinado e roteado por
    adjacência de camada. Em uma chamada isolada, todo conteúdo — inclusive
    conteúdo hostil — entra indiscriminadamente no mesmo canal do \emph{prompt},
    sem fronteira entre dado e instrução.
  \item \textbf{Defesa ativa contra manipulação do próprio agente.} Os mecanismos
    M1--M5 monitoram, detectam e contêm o comprometimento dos agentes em tempo de
    execução. Uma chamada isolada não dispõe de monitoramento comportamental nem
    de quarentena; se for envenenada ou injetada, ela simplesmente produz a saída
    manipulada.
  \item \textbf{Governança deliberativa por quórum.} As decisões estratégicas
    emergem de um colegiado que vota e registra, e não de um oráculo único. Essa
    deliberação é o que confere ao sistema robustez a participantes defeituosos,
    propriedade ausente por construção em uma chamada isolada.
\end{enumerate}

Em síntese, o MAS-Hunt não é um classificador de eventos com um \emph{prompt}
melhor; é um sistema de governança que organiza, valida e contém o comportamento
de múltiplos agentes ao redor da tarefa de caça. As condições experimentais
materializam exatamente essa distinção: C2-naive aproxima-se de um conjunto de
agentes sem governança — o regime que esta seção contrasta —, enquanto C1-full
aplica toda a pilha de mediação e contenção. A comparação entre essas condições, e
a ablação intermediária C1-nohardening, é o instrumento empírico pelo qual a
dissertação isola e mensura o valor da governança multiagente, objeto dos
Capítulos~\ref{cap:metodologia} e seguintes.
